\documentclass[11pt]{article}
\usepackage[review]{acl}
\usepackage{times}
\usepackage{latexsym}
\usepackage[T1]{fontenc}
\usepackage[utf8]{inputenc}
\usepackage{microtype}
\usepackage{inconsolata}
\usepackage{graphicx}
\usepackage{booktabs}

\title{Formality-Based Human vs. Machine Text Detection}
\author{
  Your Name \\
  CSE 188 \\
  \texttt{your-email@example.com}
}

\begin{document}
\maketitle

\begin{abstract}
This report studies whether lightweight formality features can distinguish
human-written text from machine-generated text. We train a logistic regression
classifier on features such as contraction rate, informal word rate, sentence
length statistics, punctuation usage, and capitalization. We evaluate with
Accuracy, F1, and AUROC, and include ablation and error analysis to measure
technical soundness.
\end{abstract}

\section{Introduction}

Large language models (LLMs) produce fluent text, which makes automatic source
detection increasingly important. In this project, we test a lightweight
approach: formality-based features with logistic regression.

Our hypothesis is that machine-generated text is often more stylistically
regular, while human text contains more informal variation.

Main contributions:
\begin{itemize}
  \item A simple and interpretable feature set for human-vs-machine detection.
  \item A baseline classifier with standard metrics (Accuracy, F1, AUROC).
  \item A draft analysis plan with ablation and error analysis.
\end{itemize}

\section{Related Work}

AI-generated text detection includes supervised classifiers, unsupervised
methods, and watermark-based approaches \citep{pu2022deepfake,mitchell2023detectgpt}.
Recent work also emphasizes robustness in open-domain settings \citep{li2024mage}.

Compared to large detector models, this project uses a compact and interpretable
feature-based pipeline.

\section{Method}

\subsection{Features}

We use ten formality-related features:
\begin{itemize}
  \item contraction rate
  \item informal word rate
  \item sentence length mean
  \item sentence length standard deviation
  \item exclamation rate
  \item question mark rate
  \item all-caps ratio
  \item average word length
  \item number of sentences
  \item number of words
\end{itemize}

\subsection{Model}

We train logistic regression with standardized features. This model is simple,
fast, and easy to interpret.

\section{Experimental Setup}

\subsection{Data}

We use the \texttt{ahmadreza13/human-vs-Ai-generated-dataset} setup from the
project code. The default run uses 400 samples per class (800 total), with a
stratified train/validation/test split.

\subsection{Metrics}

We report:
\begin{itemize}
  \item Accuracy
  \item F1 score (binary)
  \item AUROC
\end{itemize}

\subsection{Implementation}

The implementation uses scikit-learn logistic regression with feature scaling.
Default run command:
\begin{verbatim}
python run_train.py
\end{verbatim}

\section{Results}

Table~\ref{tab:main-results} is a placeholder for final metrics from the
training script.

\begin{table}[t]
\centering
\begin{tabular}{lccc}
\toprule
Setting & Accuracy & F1 & AUROC \\
\midrule
Default run (400/class) & TODO & TODO & TODO \\
\bottomrule
\end{tabular}
\caption{Main classification results.}
\label{tab:main-results}
\end{table}

Brief interpretation goes here:
\begin{itemize}
  \item Which metric is strongest?
  \item Where does the model fail?
  \item Is performance stable across seeds?
\end{itemize}

\section{Analysis}

\subsection{Ablation Study}

Remove one feature group at a time and measure metric drop.

\begin{table}[t]
\centering
\begin{tabular}{lccc}
\toprule
Variant & Accuracy & F1 & AUROC \\
\midrule
Full feature set & TODO & TODO & TODO \\
- contraction feature & TODO & TODO & TODO \\
- informal word feature & TODO & TODO & TODO \\
- punctuation features & TODO & TODO & TODO \\
\bottomrule
\end{tabular}
\caption{Ablation results for feature groups.}
\label{tab:ablation}
\end{table}

\subsection{Error Analysis}

Add 5--10 representative failures and explain patterns:
\begin{itemize}
  \item False positives (human predicted as machine)
  \item False negatives (machine predicted as human)
  \item Common linguistic cues behind errors
\end{itemize}

\section{Limitations}

This study has several limits:
\begin{itemize}
  \item A lightweight feature model may miss deeper semantic cues.
  \item Current experiments use a limited dataset slice and setup.
  \item Cross-domain and adversarial robustness are not fully tested yet.
\end{itemize}

Future work can add broader datasets, stronger baselines, and robustness tests.

\section{Conclusion}

This report presents a simple and interpretable baseline for human-vs-machine
text detection using formality features and logistic regression. Initial results
suggest that lightweight stylistic signals can be useful. The next step is to
improve soundness through stronger baselines, ablations, and detailed error
analysis.


\section*{Limitations}
This draft is based on a small baseline setup. Broader domains and stronger
baselines may change conclusions.

\section*{Ethics Statement}
This work focuses on text source classification and should be used carefully to
avoid over-claiming certainty in high-stakes settings.

\bibliography{custom}
\bibliographystyle{acl_natbib}

\end{document}
